\scanpage{43}
\uline{NB} Si $X_0$ est un schéma abélien \deleted{et} (ou seulement
\margin{?}\unsure{paraabélien|parabélien|pro-abélien}), alors ces conditions \deleted{impliquent} \inserted{équivalent à celle (0)} que
$\unsure{\mathfrak X|\mathcal X|\widehat X}$ soit algébrique, et dans l'énoncé des conditions
(i) à (iv) [dont alors seules (i) et (ii) sont intéressantes]
on peut supprimer la condition de projectivité.

\uline{Remarque.} Soit $S'\to S$ un \inserted{morphisme propre qui est \unsure{un}} \unsure{épimorphisme|(épimorphisme)}
effectif ($S$ loc noeth). C'est donc un morphisme
de descente effective pour les schémas relatifs plats.
Soit $X'$ plat et propre sur $S'$, muni d'une donnée de
descente à $S$, on \deleted{\unsure{veut|on}} se demande si donnée
de descente est effective.

\deleted{a) Condition locale sur $S$}

a) On peut se réduire aux localisés \inserted{$\mathcal O_{S,s}$} de $S$.

b) Moyennant une question de descente \deleted{\unsure{fi|fid}} fid.
plate, on est ramené au pb sur $\widehat{\mathcal O}_{S,s}$. Il
peut être intéressant d'avoir \deleted{des \unsure{pts}} assez de p\textsuperscript{ts}
de $X_s$, et pour ceci de faire encore une ext. finie
\deleted{\unsure{étale}} \inserted{\deleted{\unsure{étale}} $R$} de $\widehat{\mathcal O}_{S,s}$. On peut alors supposer que
\deleted{$X_s$} sur le corps résiduel $k$ de $R$, on peut
descendre, donc qu'on a un $X_0$.

% slanted
\margin{\uline{Conditions \unsure{peut|pour}} \unsure{être} \uline{\unsure{inutiles|constantes}}}
c) Supposons \inserted{la réduction précédente faite, et} que les fibres de $X'/S'$ soient
sans aut. infinitésimaux, ou que $X'$ \deleted{soit lisse sur $S'$}
ait \deleted{\illegible{1}} ``assez'' de sections \deleted{\unsure{de plus}} \deleted{\unsure{lisses sur lesquelles}}
stables par descente, le long desquelles $X'/S'$ soit
lisse. Alors \deleted{\unsure{en|par}} un argument de modules
formels permet de construire un schéma formel
$\widehat X$ sur $\widehat S$ qui descend $\widehat X'$.
% archive sheet number: 42
